\section{Selected complete-history returns and threshold completion}\label{sec:return-tools}
The first result isolates the regularity needed to turn a selected event into a moving
\emph{complete-history} return.  This distinction matters because a
fixed-time RFDE solution map may be smooth in its initial history while its
joint dependence on time and an arbitrary continuous history is not yet
smooth: before the translation part has cleared the initial interval,
changing time moves an evaluation point in a merely continuous function.

Let
\begin{equation}
 \mathcal X_{\tau_*}=C([-\tau_*,0],\R^d),\qquad
 \dot x(t)=\mathcal F(x_t),
 \qquad x_t(\theta)=x(t+\theta),                      \label{eq:general-rfde}
\end{equation}
where \(d<\infty\), \(\tau_*>0\), and
\(\mathcal F\colon U\subset\mathcal X_{\tau_*}\to\R^d\) is \(C^r\).
Write \(\Phi_t(\phi)=x_t(\phi)\).

\begin{theorem}[Eventually smooth selected-event return]
\label{thm:eventual-selected-return}
Let \(1\le k\le r\), let \(\mathcal M\) be a Banach space,
\(D\subset\mathcal M\) open, and let
\(\iota\colon D\to U\subset\mathcal X_{\tau_*}\) be \(C^k\).  Suppose all
solutions from \(\iota(D)\) exist in one common \(U\)-tube through \(T_+\).
Let \(V\subset\mathcal X_{\tau_*}\) be open, let
\(g\colon V\to\R\) be \(C^k\), and assume
\(\Phi_t(\iota(u))\in V\) for \((t,u)\in I\times D\).  Put
\[
 \mathcal H(t,u)=g(\Phi_t(\iota(u))).
\]
Suppose a common interval \(I=[T_-,T_+]\) satisfies
\begin{equation}
 T_->k\tau_*,\qquad
 \sup_{u\in D}\mathcal H(T_-,u)\le-\delta_-,\qquad
 \inf_{u\in D}\mathcal H(T_+,u)\ge\delta_+,
 \qquad
 \partial_t\mathcal H\ge a_*>0                     \label{eq:eventual-return-gates}
\end{equation}
on \(I\times D\), for \(\delta_-,\delta_+,a_*>0\).  Then each
\(u\in D\) has exactly one selected event
\(T(u)\in(T_-,T_+)\), and the event time and ambient event-hit map
\[
 T\colon D\to\R,
 \qquad \mathcal E(u)=\Phi_{T(u)}(\iota(u))
       \colon D\to\mathcal X_{\tau_*}
\]
are \(C^k\).  The threshold \(k\tau_*\) is a uniform sufficient condition;
no necessity or optimality is asserted.

If \(\iota\) is a \(C^k\) chart of a local section \(\Sigma\), and
\(\mathcal E(D_0)\subset\iota(D)\) for an open \(D_0\subset D\), then
\begin{equation}
 \mathcal Q=\iota^{-1}\circ\mathcal E\big|_{D_0}
 \colon D_0\to D                                      \label{eq:induced-selected-return}
\end{equation}
is the induced \(C^k\) selected section-return map.
\end{theorem}

\begin{proof}
On the common tube, repeated Fr\'echet differentiation of the Volterra
equation gives, for every fixed time, the triangular initial-data
variational hierarchy through order \(k\).  We record the additional joint
smoothing step because it is essential here.  Inductively in \(j\), the
mixed jets
\[
 \partial_t^aD_\phi^b x(t;\phi),
 \qquad a+b\le j,
\]
exist after \((j-1)\tau_*\) and depend continuously on \((t,\phi)\) in the
corresponding multilinear-operator norms.  For \(j=1\) this follows from the
integral equation and the first variational equation.  If it holds through
order \(j\), differentiating \(x'=\mathcal F(x_t)\) and the triangular
variational equations once more expresses every order-\(j+1\) jet as a
finite sum of continuous derivatives of \(\mathcal F\) applied to lower-order
history jets.  After one further maximum delay, every argument segment lies
in the already regularized interval.  The linear inhomogeneous variational
equations and Gronwall's inequality give continuity in the multilinear-
operator norms, completing the induction.

The earliest physical time in \(\Phi_t(\phi)\) is \(t-\tau_*\).  Hence, for
every total order \(j\le k\), the preceding jets define continuous
\(\mathcal X_{\tau_*}\)-valued derivatives once \(t>j\tau_*\).  Therefore
\((t,\phi)\mapsto\Phi_t(\phi)\) is jointly \(C^k\) for \(t>k\tau_*\), and
\(\mathcal H\) is jointly \(C^k\) on \(I\times D\).  The endpoint signs and
positive speed give existence and uniqueness in the declared window; the
Banach implicit-function theorem gives \(T\in C^k\), and composition gives
\(\mathcal E\in C^k\).  The last assertion follows by composing with the
section chart inverse.  See also the fixed-time RFDE and Poincar\'e-map
background in \citet[Chs.~VII.4 and XIV.3]{diekmann1995delay}.
\end{proof}

We now state the exact implication that would turn the remaining model
estimates into a physical threshold.  It is deliberately conditional: the
hypotheses involving the effective stable graph and the two attraction tubes
have not yet been verified for \cref{eq:scalar-v,eq:scalar-w}.

Let \(\xi=(a,\kappa_3)\) range in a closed rectangle
\(\Xi\subset\R^2\) about \(\xi_0=(1/4,1/200)\), with the other parameters
fixed as in \cref{eq:center}.  Denote the quiet equilibrium, inner and outer
periodic orbits, pulse release history, and event-aligned section entrance by
\(E_q(\xi)\), \(\Gamma_i(\xi)\), \(\Gamma_o(\xi)\), \(R_\xi(J)\), and
\(B_\xi(J)\), respectively.  In section coordinates
\((u,z)\in E^u_\xi\oplus E^s_\xi\), write
\begin{equation}
 W^s_{\rm loc}(\Gamma_i(\xi))\cap\Sigma_\xi
 =\{(u,z):u=\psi_\xi(z)\},
 \qquad
 H(\xi,J)=u_\xi(J)-\psi_\xi(z_\xi(J)).                 \label{eq:parametric-pulse-gap}
\end{equation}

\begin{theorem}[Conditional pulse-threshold completion]
\label{thm:pulse-completion}
Assume the following objects and estimates hold with common domains, radii,
and strict constants for every \(\xi\in\Xi\).  The equilibrium
\(E_q(\xi)\) and phase-fixed branches \(\Gamma_i(\xi),\Gamma_o(\xi)\) are
\(C^1\) in \(\xi\), and the quiet equilibrium and outer cycle possess
quantitative attraction tubes.  The inner section return has a
one-dimensional unstable bundle and a stable complement, both \(C^1\) in
\(\xi\), and \(\psi_\xi\) is jointly \(C^1\) in \((\xi,z)\) on one common
graph domain.

Assume also that \((\xi,J)\mapsto R_\xi(J)\) and the event branch
\(B_\xi(J)\) are jointly \(C^1\) on
\(\Xi\times I_J\), where \(I_J=[J_-,J_+]\).  The event does not switch
branches, its speed is bounded away from zero, and every stable coordinate
lies in the common domain of \(\psi_\xi\).  Suppose, after fixing the
orientation, that there are \(\eta_H,m_J>0\) such that
\begin{equation}
 H(\xi,J_-)\ge\eta_H,
 \qquad H(\xi,J_+)\le-\eta_H,
 \qquad \partial_JH(\xi,J)\le-m_J<0                   \label{eq:gap-crossing}
\end{equation}
on the whole product.  Finally, suppose the signed-exit dynamics sends the
\(H>0\) side into a complete-history exit slab attached to the quiet
attraction tube and the \(H<0\) side into one attached to the outer
attraction tube, uniformly in \(\xi\), with no other local exit; histories on
\(H=0\) lie on \(W^s_{\rm loc}(\Gamma_i(\xi))\).

Then there is a unique \(C^1\) function
\(J_c:\Xi\to(J_-,J_+)\) satisfying \(H(\xi,J_c(\xi))=0\), with
\begin{equation}
 D_\xi J_c(\xi)
 =-\frac{D_\xi H(\xi,J_c(\xi))}
          {\partial_JH(\xi,J_c(\xi))}.                 \label{eq:threshold-derivative}
\end{equation}
Moreover, in the complete history space,
\begin{equation}
 \begin{aligned}
 J<J_c(\xi)&\Longrightarrow
   \norm{\Phi_t^\xi R_\xi(J)-E_q(\xi)}_X\longrightarrow0,\\
 J=J_c(\xi)&\Longrightarrow
   \operatorname{dist}_X(\Phi_t^\xi R_\xi(J),\Gamma_i(\xi))
      \longrightarrow0,\\
 J>J_c(\xi)&\Longrightarrow
   \operatorname{dist}_X(\Phi_t^\xi R_\xi(J),\Gamma_o(\xi))
      \longrightarrow0
 \end{aligned}
 \qquad(t\to\infty).                                  \label{eq:routing}
\end{equation}
\end{theorem}

\begin{proof}
The endpoint signs and strict derivative give one and only one zero for each
\(\xi\).  The implicit-function theorem yields the stated parameter
dependence and derivative.  The stable graph identifies the zero with the
local stable sheet, while the signed-exit and attraction-tube assumptions give
the three asymptotic alternatives in \cref{eq:routing}.
\end{proof}

For the concrete center parameter, the event branch and several unadjusted
functional signs are proved below.  The common quantitative graph,
graph-adjusted endpoint signs, selected crossing, and routing assumptions of
\cref{thm:pulse-completion} remain open.

No exclusion of earlier hits is used in
\cref{thm:eventual-selected-return}: it is needed only to call the branch a
first return or to assign it an exact ordinal.  This separation also avoids
an unnecessary algebraic restriction on a stable-manifold construction.

\begin{lemma}[Direct selected return identifies the stable-set germ]
\label{lem:direct-selected-stable-germ}
Let \(\Phi\) be a continuous semiflow, \(\Gamma\) a periodic orbit through a
transverse local section \(\Sigma\) at \(p\).  Choose a local section patch
\(N\subset\Sigma\) with the phase-isolation property
\begin{equation}
 x_n\in N,\qquad \operatorname{dist}(x_n,\Gamma)\to0
 \quad\Longrightarrow\quad x_n\to p.                 \label{eq:section-phase-isolation}
\end{equation}
Let \(\mathcal Q(x)=\Phi_{\Theta(x)}(x)\) be a selected local return of the
same semiflow with
\[
 \mathcal Q\colon N\to N,\qquad \mathcal Q(p)=p,\qquad
 \Theta(p)=mP,\qquad
 0<\Theta_-\le\Theta(x)\le\Theta_+<\infty.
\]
Suppose every intervening arc between successive iterates stays in one common
flow tube.  Restrict both stable sets to points whose relevant iterates or
flow arcs stay in these declared local domains, and assume that local
flow-stable points in \(N\) have the selected hits recurrently defined in
\(N\).  Then, as germs at \(p\),
\begin{equation}
 \operatorname{germ}_p W^s_N(\mathcal Q)
 =\operatorname{germ}_p\bigl(N\cap W^s_{\rm tube}(\Gamma)\bigr).
                                                               \label{eq:direct-return-stable-germ}
\end{equation}
Consequently, a direct near-\(mP\) selected return can construct the intrinsic
periodic-orbit stable-sheet germ; an identity \(\mathcal Q=P^m\) and a
first-return label are sufficient conveniences, not necessary hypotheses.
If \(\mathcal Q\) is \(C^2\), \(p\) has the required hyperbolic splitting,
and a local graph theorem applies, its stable graph represents this same
intrinsic germ.
\end{lemma}

\begin{proof}
Put
\(t_n=\sum_{j=0}^{n-1}\Theta(\mathcal Q^j x)\).  If
\(\mathcal Q^n(x)\to p\), then \(t_n\to\infty\) by the positive lower bound;
joint continuity of \(\Phi_t\), uniformly for
\(0\le t\le\Theta_+\), sends every intervening late-time arc to the
corresponding arc of \(\Gamma\).  Conversely, convergence of the flow orbit
to \(\Gamma\), together with recurrent selected hits in \(N\), gives
\(\operatorname{dist}(\mathcal Q^n x,\Gamma)\to0\).  The phase-isolation
property \cref{eq:section-phase-isolation} then gives
\(\mathcal Q^n x\to p\).
\end{proof}

